\input{../tex-inputs/latex-leseansicht-vorspann}

\section[Gustav Schwarzkopf an Arthur Schnitzler, 26. 1. 1914]{L04254 Gustav Schwarzkopf an Arthur Schnitzler, 26. 1. 1914}
\nopagebreak\mylabel{L04254v}
\rehead{ }\normalsize\beginnumbering\briefempfaengerindex{Schnitzler, Arthur@\textsc{Schnitzler, Arthur}!zzzSchwarzkopf, Gustav@\emph{von Gustav Schwarzkopf}!1914-01-261@{26. 1. 1914}|(be}
\toendnotes[C]{\smallbreak\pagebreak[2]}
\correspDesc{Versand  durch Gustav Schwarzkopf am 26. 1. 1914 in Wien
\newline{}Erhalt  durch Arthur Schnitzler im Zeitraum [27. 1. 1914 – 31. 1. 1914?] in München}\toendnotes[C]{\smallbreak}
\Standort{CUL, Schnitzler, B 96a.}
\physDesc{Postkarte, 521 Zeichen
\newline{}Handschrift: schwarze Tinte, deutsche Kurrent
\newline{}Versand: Stempel: »\nobreak{}\oindex{I., Innere Stadt@\textbf{I., Innere Stadt}, \emph{Verwaltungsgebiet}|pwk}1/1 Wien 8, 26. I. 14, 11\nobreak{}«.  }\toendnotes[C]{\smallbreak}\pstart{}{\pb}Herrn\pend{}\pstart{}\textsc{D\textsuperscript{r} Arthur Schnitzler}\pend{}\pstart{}\textsc{München\oindex{München@\textbf{München}|pw}}\pend{}\pstart{}\textsc{Hotel Marienbad\oindex{Hotel Marienbad [München]@\textbf{Hotel Marienbad [München]}, \emph{Hotel}|pw}}\pend{}{\bigskip}\vspace{1em}
\pstart
           \raggedleft{}{\pb}26. I. 14\pend
           \vspace{0.5em}
\pstart
           Lieber Arthur, ich bin{ }ſchon wieder gerettet, aber die anhaltende
               große Kälte bringt mich noch zu Hauſe zu bleiben. Ka{\geminationn}
               alſo auch heute nicht zur \label{K_L04254-1v}\edtext{Generalprobe der Liebelei\pwindex{\textcolor{red}{\textsuperscript{XXXX indx1}}!Elskovsleg@\strich\emph{Elskovsleg}|pw}\pwindex{\textcolor{red}{\textsuperscript{XXXX indx1}}!Elskovsleg@\strich\emph{Elskovsleg}|pw}\eventindex{Elitekino@\textbf{Elitekino}!Generalprobe von Liebelei [Elskovsleg], 26.1.1914@Generalprobe von Liebelei [Elskovsleg], 26.1.1914|pwv}}{\lemma{\textnormal{\emph{Generalprobe der Liebelei}}}\Cendnote{\textnormal{Es handelte sich um ein Vorabscreening\eventindex{Elitekino@\textbf{Elitekino}!Generalprobe von Liebelei [Elskovsleg], 26.1.1914@Generalprobe von Liebelei [Elskovsleg], 26.1.1914|pwkv} des dänischen\oindex{Dänemark@\textbf{Dänemark}|pwk} Films
                  \emph{Elskovsleg}\pwindex{\textcolor{red}{\textsuperscript{XXXX indx1}}!Elskovsleg@\strich\emph{Elskovsleg}|pwk}\pwindex{\textcolor{red}{\textsuperscript{XXXX indx1}}!Elskovsleg@\strich\emph{Elskovsleg}|pwk} mit deutschen Zwischentiteln, basierend auf Schnitzlers{ }\emph{Liebelei}\pwindex{Schnitzler, Arthur 15.\,5.\,1862 Wien – 21.\,10.\,1931 ebd.@\textsc{Schnitzler, Arthur} (15.\,5.\,1862 Wien – 21.\,10.\,1931 ebd.), \emph{Schriftsteller, Mediziner}!Liebelei. Schauspiel in drei Akten@\strich\emph{Liebelei. Schauspiel in drei Akten}|pwk}. Es fand im Elitekino\oindex{Wien@\textbf{Wien}!I., Innere Stadt@\textbf{I., Innere Stadt}!Elitekino@\textbf{Elitekino}, \emph{Kino}|pwk}
                  statt.}}}\label{K_L04254-1} gehen — besten Dank für die Karte. Die \label{K_L04254-2v}\edtext{drei Einakter\pwindex{Wedekind, Frank 24.\,7.\,1864 Hannover – 9.\,3.\,1918 München@\textsc{Wedekind, Frank} (24.\,7.\,1864 Hannover – 9.\,3.\,1918 München), \emph{Schriftsteller, Schauspieler, Schriftsteller}!Kammersänger@\strich\emph{Der Kammersänger}|pwv}\pwindex{Courteline, Georges 25.\,6.\,1858 Tours – 25.\,6.\,1929 Paris@\textsc{Courteline, Georges} (25.\,6.\,1858 Tours – 25.\,6.\,1929 Paris), \emph{Schriftsteller}!Boubouroche. Lebensbild in 2 Acten@\strich\emph{Boubouroche. Lebensbild in 2 Acten}|pwv}\pwindex{Schnitzler, Arthur 15.\,5.\,1862 Wien – 21.\,10.\,1931 ebd.@\textsc{Schnitzler, Arthur} (15.\,5.\,1862 Wien – 21.\,10.\,1931 ebd.), \emph{Schriftsteller, Mediziner}!Literatur@\strich\emph{Literatur}|pwv}{ }ſind für 31.{ }angeſetzt\eventindex{Burgtheater@\textbf{Burgtheater}!Premiere von Der Kammersänger, Boubouroche, Literatur, 30.1.1914@Premiere von Der Kammersänger, Boubouroche, Literatur, 30.1.1914|pwv}}{\lemma{\textnormal{\emph{drei … angesetzt}}}\Cendnote{\textnormal{Tatsächlich fand am 30. 1. 1914
                  am \emph{Burgtheater}\orgindex{Burgtheater@Burgtheater|pwk} die Premiere von \emph{Der Kammersänger}\pwindex{Wedekind, Frank 24.\,7.\,1864 Hannover – 9.\,3.\,1918 München@\textsc{Wedekind, Frank} (24.\,7.\,1864 Hannover – 9.\,3.\,1918 München), \emph{Schriftsteller, Schauspieler, Schriftsteller}!Kammersänger@\strich\emph{Der Kammersänger}|pwk} (Frank Wedekind\pwindex{Wedekind, Frank 24.\,7.\,1864 Hannover – 9.\,3.\,1918 München@\textsc{Wedekind, Frank} (24.\,7.\,1864 Hannover – 9.\,3.\,1918 München), \emph{Schriftsteller, Schauspieler, Schriftsteller}|pwk}), \emph{Boubouroche. Lebensbild in 2 Acten}\pwindex{Courteline, Georges 25.\,6.\,1858 Tours – 25.\,6.\,1929 Paris@\textsc{Courteline, Georges} (25.\,6.\,1858 Tours – 25.\,6.\,1929 Paris), \emph{Schriftsteller}!Boubouroche. Lebensbild in 2 Acten@\strich\emph{Boubouroche. Lebensbild in 2 Acten}|pwk} (Georges Courteline\pwindex{Courteline, Georges 25.\,6.\,1858 Tours – 25.\,6.\,1929 Paris@\textsc{Courteline, Georges} (25.\,6.\,1858 Tours – 25.\,6.\,1929 Paris), \emph{Schriftsteller}|pwk}) und
                      \emph{Literatur}\pwindex{Schnitzler, Arthur 15.\,5.\,1862 Wien – 21.\,10.\,1931 ebd.@\textsc{Schnitzler, Arthur} (15.\,5.\,1862 Wien – 21.\,10.\,1931 ebd.), \emph{Schriftsteller, Mediziner}!Literatur@\strich\emph{Literatur}|pwk} (Schnitzler)\eventindex{Burgtheater@\textbf{Burgtheater}!Premiere von Der Kammersänger, Boubouroche, Literatur, 30.1.1914@Premiere von Der Kammersänger, Boubouroche, Literatur, 30.1.1914|pwk} statt. Schnitzler
                  besuchte erst die Vorstellung\eventindex{Burgtheater@\textbf{Burgtheater}!Aufführung von Der Kammersänger, Boubouroche, Literatur, 6.2.1914@Aufführung von Der Kammersänger, Boubouroche, Literatur, 6.2.1914|pwkv} am 6. 2. 1914.
               }}}\label{K_L04254-2}, – nur das wiſſen Sie wol{ }ſchon. We{\geminationn} Sie auf Berlin\oindex{Berlin@\textbf{Berlin}, \emph{Hauptstadt}|pw} verzichten, werden Sie wol{ }ſchon hier{ }ſein und mir wird hoffentlich
               das Wetter erlauben auch dabei zu{ }ſein.\pend
           
\pstart
           Ihre Frau\pwindex{Schnitzler, Olga 17.\,1.\,1882 Wien – 13.\,1.\,1970 Lugano@\textsc{Schnitzler, Olga} (17.\,1.\,1882 Wien – 13.\,1.\,1970 Lugano), \emph{Schauspielerin, Sängerin}|pwv} und Sie herzlichſt grüßend{\\[\baselineskip]}Ihr{\\[\baselineskip]}\spacefill\mbox{Gustav}\pend
           \leftskip=0em{}
\pstart
           \noindent{}Bitte, grüßen Sie \textsc{Paul Marx\pwindex{Marx, Paul 21.\,7.\,1879 Wien – 30.\,10.\,1956 ebd.@\textsc{Marx, Paul} (21.\,7.\,1879 Wien – 30.\,10.\,1956 ebd.), \emph{Regisseur, Schauspieler}|pw}}.\pend
           \selectlanguage{ngerman}\endnumbering\briefempfaengerindex{Schnitzler, Arthur@\textsc{Schnitzler, Arthur}!zzzSchwarzkopf, Gustav@\emph{von Gustav Schwarzkopf}!1914-01-261@{26. 1. 1914}|)be}\mylabel{L04254h}
\begin{anhang}
\end{anhang}\newcommand{\dateiname}{L04254}\newcommand{\titel}{Gustav Schwarzkopf an Arthur Schnitzler, 26. 1. 1914}\newcommand{\editorInnen}{Herausgegeben von Jahnke, SelmaMüller, Martin Anton}\input{../tex-inputs/latex-leseansicht-abspann}
